\documentclass[11pt, a4paper]{ctexart}

% 页面与边距设置
\usepackage{geometry}
\geometry{left=2.5cm, right=2.5cm, top=3.0cm, bottom=3.0cm}

% 常用宏包
\usepackage{amsmath, amssymb}
\usepackage{enumitem}
\usepackage{xcolor}
\usepackage{titlesec}
\usepackage{listings}

% 颜色与代码块样式定义
\definecolor{codegreen}{rgb}{0,0.6,0}
\definecolor{codegray}{rgb}{0.5,0.5,0.5}
\definecolor{codepurple}{rgb}{0.58,0,0.82}
\definecolor{backcolour}{rgb}{0.96,0.96,0.94}

\lstdefinestyle{mystyle}{
    backgroundcolor=\color{backcolour},   
    commentstyle=\color{codegreen},
    keywordstyle=\color{blue},
    numberstyle=\tiny\color{codegray},
    stringstyle=\color{codepurple},
    basicstyle=\ttfamily\small,
    breakatwhitespace=false,         
    breaklines=true,                 
    captionpos=b,                    
    keepspaces=true,                 
    numbers=left,                    
    numbersep=5pt,                  
    showspaces=false,                
    showstringspaces=false,
    showtabs=false,                  
    tabsize=4
}
\lstset{style=mystyle}

\setlist{nosep, itemsep=1ex, topsep=1ex}

\begin{document}

\begin{center}
    {\LARGE \textbf{面向边缘计算的 MD-SLAM 系统：\\ [0.5em] 核心方法与实现细节}} \\ [2em]
\end{center}

\section{创新点 1：高频视觉先验与曼哈顿结构豁免}

\subsection{详细方法与流程}
针对边缘计算设备（如限制算力的移动端）在高度动态场景下的实时性衰减与特征误杀问题，本系统设计了基于双线程架构的级联概率滤波机制[cite: 1]。

\begin{itemize}[label=$\bullet$]
    \item \textbf{高频视觉先验：} 在前端追踪线程中，部署边缘端优化的轻量级 YOLO 实例分割网络，利用边界框与掩码进行特征的极速初始筛查，提供高频视觉先验[cite: 1]。
    \item \textbf{级联概率滤波拒绝硬剔除：} 摒弃将动态边界框内特征全部直接“硬剔除”的传统做法，所有嫌疑特征将进入后续的级联概率滤波管线[cite: 1]。
    \item \textbf{曼哈顿结构豁免：} 在局部建图线程中，提取符合曼哈顿世界假设的正交平面[cite: 1]。若嫌疑特征点完美拟合这些正交结构（如墙壁、地板），算法将触发结构豁免机制，直接赋予其 $P_{dynamic} = 0$ 的动态概率，强制保留这些被动态物体遮挡或干扰的静态背景特征[cite: 1]。
\end{itemize}

\subsection{数学建模}
曼哈顿正交平面在三维空间中被数学公式化提取，表达为：
$$n^T P + d = 0$$[cite: 1]
定义曼哈顿结构掩码函数 $M_{manhattan}$，当特征点属于该平面时，直接绕过极线几何校验，分配动态概率：
$$P_{dynamic}(P_i | M_{manhattan}(P_i) = 1) = 0$$[cite: 1]

\subsection{算法伪代码}
\begin{lstlisting}[language=Python, caption={Visual Prior and Manhattan Immunity}]
# Input: Features {F}, Edge-YOLO Mask, Depth Map
# Output: Classified Features with Prior Probability

static_features = []
suspect_features = []

# Extract orthogonal planes based on Manhattan World assumption
manhattan_planes = extract_orthogonal_planes(depth_map) 

for pt in F:
    # High-frequency visual prior screening
    if YOLO_Mask[pt.u, pt.v] == STATIC:
        pt.p_dynamic = 0.0
        static_features.append(pt)
    else:
        # Structural Immunity Mechanism
        is_immune = False
        pt_3d = BackProject(pt.u, pt.v, depth)
        
        for plane in manhattan_planes:
            if abs(dot(plane.n, pt_3d) + plane.d) < THRESHOLD:
                is_immune = True
                break
                
        if is_immune:
            pt.p_dynamic = 0.0 # Bypassing epipolar checks
            static_features.append(pt)
        else:
            suspect_features.append(pt)
            
return static_features, suspect_features
\end{lstlisting}

\vspace{1em}

\section{创新点 2：极线几何校验与平面阴影剔除}

\subsection{详细方法与流程}
对于未获得曼哈顿结构豁免的其余特征点，系统将执行几何与光度层面的深度校验，以解决间接动态干扰（如平滑表面上的反射与阴影）造成的轨迹漂移[cite: 1]。

\begin{itemize}[label=$\bullet$]
    \item \textbf{极线几何校验：} 针对非曼哈顿平面上的点，计算其极线几何误差（桑普森距离），以此推导得出基于几何特征的动态概率 $P_{geo}$[cite: 1]。
    \item \textbf{平面阴影剔除：} 动态物体产生的移动阴影极易被误认为静态角点。系统将潜在的阴影区域通过数学公式投影至已提取的曼哈顿地面（Ground Plane）上，并为该区域内的特征点分配阴影先验概率，执行精准剔除[cite: 1]。
\end{itemize}

\subsection{数学建模}
极线几何校验依赖计算桑普森距离（Sampson distance）来推导 $P_{geo}$[cite: 1]。同时，平面阴影剔除模型将阴影先验概率 $P_{shadow}$ 与 $P_{geo}$ 融合，更新该点的最终动态概率 $P_{dynamic}$。

\subsection{算法伪代码}
\begin{lstlisting}[language=Python, caption={Epipolar Verification and Shadow Rejection}]
# Input: Suspect Features {S}, Fundamental Matrix F, Ground Plane
# Output: Features with evaluated P_dynamic

valid_features = []
shadow_regions = project_shadows_to_plane(dynamic_objects, ground_plane)

for pt in S:
    # 1. Epipolar Geometry Verification
    sampson_dist = compute_sampson_distance(pt, F)
    p_geo = derive_geometric_probability(sampson_dist) 
    
    # 2. Planar Shadow Rejection
    p_shadow = 0.0
    if is_inside_projected_shadow(pt, shadow_regions):
        p_shadow = assign_prior_shadow_probability() 
        
    # Fuse probabilities
    pt.p_dynamic = max(p_geo, p_shadow)
    
    if pt.p_dynamic < DYNAMIC_THRESHOLD:
        valid_features.append(pt)

return valid_features
\end{lstlisting}

\vspace{1em}

\section{创新点 3：密度感知的自适应点线联合图优化}

\subsection{详细方法与流程}
在极端纹理退化或动态目标过大导致特征被大规模剔除的场景下，系统容易面临特征匮乏（Feature Starvation）与跟踪丢失[cite: 1]。后端设计了密度感知的自适应点线联合 Bundle Adjustment (BA) 来解决此问题[cite: 1]。

\begin{itemize}[label=$\bullet$]
    \item \textbf{特征密度评估：} 系统实时计算当前参与优化的活跃特征数量 $N_{active}$，并定义特征密度惩罚系数 $\gamma$[cite: 1]。
    \item \textbf{自适应信息矩阵缩放：} 基于贝叶斯动态概率，对优化图中的信息矩阵（Information Matrices）进行动态自适应缩放，抑制高概率动态点的权重，防止轨迹漂移[cite: 1]。
    \item \textbf{点线联合约束：} 当点特征出现严重退化或被施加高额惩罚时，系统自动引入曼哈顿线特征（Line Features）提供鲁棒的方向约束，实现联合优化兜底[cite: 1]。
\end{itemize}

\subsection{数学建模}
活跃特征数量记为 $N_{active}$，对应密度惩罚系数为 $\gamma$[cite: 1]。
后端自适应信息矩阵缩放的核心公式为：
$$\Omega_{i}^{new} = W_{s} \cdot \Omega_{i}$$[cite: 1]
其中，权重因子 $W_s$ 由动态概率直接推导得出：
$$W_{s} = 1 - P_{dynamic}$$[cite: 1]

\subsection{算法伪代码}
\begin{lstlisting}[language=C++, caption={Density-Aware Adaptive Point-Line BA}]
// Input: Valid Points {P}, Manhattan Lines {L}, Pose T
// Output: Optimized Pose T_opt

void DensityAwareOptimization(std::vector<MapPoint>& points, Matrix4d& T) {
    Optimizer optimizer;
    optimizer.addPoseVertex(T);
    
    // 1. Feature Density Assessment
    int N_active = calculate_active_features(points); 
    double gamma = compute_density_penalty_coefficient(N_active); 
    
    // 2. Adaptive Information Matrix Scaling
    for (auto& pt : points) {
        auto edge = new PointReprojectionEdge(pt);
        
        // W_s = 1 - P_dynamic
        double W_s = 1.0 - pt.p_dynamic; 
        
        // Omega_new = W_s * Omega_i
        Matrix2d Omega_new = W_s * IdentityMatrix(); 
        edge->setInformation(Omega_new);
        optimizer.addEdge(edge);
    }

    // 3. Point-Line Joint Constraint Fallback
    if (N_active < STARVATION_THRESHOLD) {
        for (auto& line : manhattan_lines) {
            auto edge_line = new LineProjectionEdge(line);
            // Robust orientation constraints when point features are degraded
            edge_line->setInformation(IdentityMatrix() * HIGH_WEIGHT); 
            optimizer.addEdge(edge_line);
        }
    }

    optimizer.optimize();
    T = optimizer.getOptimizedPose();
}
\end{lstlisting}

\end{document}